
!TeX root = ManuscriptGuide.tex
\newpage
\clearpage
\section{Technical Presentation}

\begin{verbatim}

1. Title
    1.1 gPCD: A GPU-Driven Parallel Collision Detection Framework
    1.2 Author
    1.3 University of Florida
    1.4 Conference / Venue

2. Motivation
    2.1 Particle-based simulation applications
    2.2 Importance of collision detection
    2.3 Computational challenges at large particle counts
    2.4 Need for scalable GPU solutions

3. Problem Statement
    3.1 Collision detection as a performance bottleneck
    3.2 Memory-access challenges
    3.3 Synchronization challenges
    3.4 Scalability requirements
    3.5 Objectives of gPCD

4. Background
    4.1 Particle collision detection workflow
    4.2 Broad-phase collision detection
    4.3 Narrow-phase collision detection
    4.4 Uniform-grid approaches
    4.5 GPU-based collision detection approaches

5. gPCD Framework Overview
    5.1 Design objectives
    5.2 Fully GPU-resident execution
    5.3 Graphics-pipeline occupancy generation
    5.4 Compute-shader collision detection
    5.5 Overall processing pipeline

6. Cell Array Organization
    6.1 Uniform-grid spatial decomposition
    6.2 Cell indexing
    6.3 Occupancy storage
    6.4 Sparse occupancy handling
    6.5 O(1) cell access
    6.6 No empty-cell processing

7. Graphics-Pipeline Broad Phase
    7.1 Particle representation
    7.2 AABB construction
    7.3 Eight-corner generation
    7.4 Vertex-shader processing
    7.5 Fragment-shader occupancy generation
    7.6 Occupancy-list creation

8. Compute-Based Narrow Phase
    8.1 Candidate-pair generation
    8.2 Exact collision testing
    8.3 Parallel execution model
    8.4 Source-owned updates
    8.5 Collision determination

9. Experimental Setup
    9.1 Hardware platform
    9.2 GPU specifications
    9.3 Test configurations
    9.4 Random particle distributions
    9.5 Verification mode
    9.6 Performance mode
    9.7 Measurement methodology

10. Scaling Results
    10.1 Execution time versus particle count
    10.2 Linear scaling behavior
    10.3 Regression analysis
    10.4 Discussion of scaling characteristics

11. Throughput Results
    11.1 Throughput definition
    11.2 Throughput versus particle count
    11.3 Peak throughput
    11.4 Largest benchmark configuration
    11.5 Performance observations

12. Time-per-Particle Analysis
    12.1 Time-per-particle definition
    12.2 Scaling efficiency
    12.3 Constant-cost behavior
    12.4 Interpretation of results

13. Memory Locality Study
    13.1 Motivation
    13.2 Sequential occupancy layout
    13.3 Random occupancy layout
    13.4 Performance comparison
    13.5 Locality effects
    13.6 Implications for GPU execution

14. Discussion
    14.1 Exact collision detection
    14.2 GPU-driven occupancy generation
    14.3 Sparse-domain advantages
    14.4 Memory-access behavior
    14.5 Practical scalability considerations

15. Contributions
    15.1 Graphics-pipeline occupancy generation
    15.2 Fully GPU-resident collision detection
    15.3 Exact collision detection
    15.4 O(1) cell access
    15.5 Near-linear scaling
    15.6 Multi-million-particle capability

16. Future Work
    16.1 Larger GPU configurations
    16.2 Multi-GPU execution
    16.3 Distributed computing
    16.4 Thermo-fluid simulation integration
    16.5 Additional large-scale studies

17. Conclusions
    17.1 Summary of gPCD
    17.2 Major findings
    17.3 Performance achievements
    17.4 Scalability observations
    17.5 Closing remarks

18. Questions
    18.1 Discussion

\end{verbatim}
